\section{Сложные случаи функционирования адаптивного РНМ и направления развития}
\label{sec:ch1_complex_cases}

В данном разделе описаны современные проблемы, возникающие с функционирования \commentnote{DISS-005-C0050}{Евдаков Алексей}{Здесь и далее актуально бы указать отрасль, в которой возникало.}{адаптивного }органа РНМ, разбор данных событий, идеи об их решении в рамках существующих параметрах и выводы о том, насколько все проблемы можно считать решаемыми в существующих рамках.

\subsection{Адаптивность по времени РНМ (колебание)}

С внедрением адаптивной функции подстройки амплитуды уставки РНМ значительно улучшилось качество функционирования БАВР. Однако, на ряде предприятий возникла проблема колебательного характера нагрузки (рис. 1). Данный фактор характерен для двигателей с частотным приводом.

\begin{figure}[htbp]
    \centering
    \includegraphics[width=0.85\textwidth]{image21.png}
    \caption[Колебание амплитуды тока]{Колебание амплитуды тока\newline{\small(Файл можно найти в папке «\_вспомогательные файлы\textbackslash{}Краснодар \commentnote{DISS-005-C0051}{Евдаков Алексей}{t211131102963\_o0309}{Водоканал})}}
    \label{fig:ch1_img21}
\end{figure}

\subsubsection{Рекомендации по изменению параметров БАВР}

Ниже приведён список рекомендаций по корректировке параметров для улучшения функционирования в данных случаях:

\begin{itemize}[leftmargin=*,itemsep=0pt]
    \item Для повышения надёжности функционирования органа РНМ -- установить уставку по времени кратно периоду колебаний.
\end{itemize}

\begin{figure}[htbp]
    \centering
    \includegraphics[width=0.7\textwidth]{image22.png}
    \label{fig:ch1_img22}
\end{figure}

\textbf{Примечание:} приведённые скриншоты могу немного отличать в разных версиях ПО.

\subsubsection{Примечание по анализу осциллограмм}

Отметим, что период и амплитуда колебаний могут отличаться для различных двигателей. Поэтому подстройка данного параметра от стандартных значений требуется на каждом объекте. К тому же, амплитуда колебаний -- зависит от характера нагрузки на устройства, и может как отличаться, так и практически отсутствовать для конкретного двигателя (таблица~1). При этом время срабатывания БАВР может получать довольно большим из-за сложностей работы БАВР. Для случаев же с стандартными реализациями РНМ -- стабильной работы практически невозможно добиться, или необходимо сильно занижать уставку.

Наиболее критичен случай, когда функционируют двигатели с разными периодами колебаний, так как уставку приходится задавать как наибольшее общее кратное, что может оказаться слишком большим значением. Автоматически же данный параметр не изменяется, да и в целом требуется очень сложной ручной подстройки.

% DISS-005-B0410 (Заголовок Таблицы 1 и замечания C0052..C0054)
\paragraph{Таблица 1. Работа системы \commentnote{DISS-005-C0052}{Евдаков Алексей}{Добавить столбец с периодом колебания}{\commentnote{DISS-005-C0053}{Евдаков Алексей}{+ схема подстанции (движки) + где замеры.}{\commentnote{DISS-005-C0054}{Евдаков Алексей}{+Добавить статистику с осциллограмм}{БАВР}}}}
\label{tab:bavr_operation_title}

\begin{table}[htbp]
\centering
\small
\begin{tabularx}{\textwidth}{|c|>{\centering\arraybackslash}X|c|c|c|c|>{\centering\arraybackslash}p{2.8cm}|}
\hline
\textbf{№ эксп.} & \textbf{\begin{tabular}[c]{@{}c@{}}№ работающего ПЧ /\\ Срабатывание БАВР\end{tabular}} & \textbf{\begin{tabular}[c]{@{}c@{}}Ввод 1\\ (Т1)\end{tabular}} & \textbf{\begin{tabular}[c]{@{}c@{}}Ввод 2\\ (Т2)\end{tabular}} & \textbf{$i_{\max}$} & \textbf{$K_{\text{броск}}$} & \textbf{\commentnote{DISS-005-C0055}{Feel}{Вот это ч хороший столбец\newline А вывод-то какой?}{\commentnote{DISS-005-C0056}{Евдаков Алексей}{Никакой)) Я прост вставил таблицу, в которой у меня были указаны эти колебания в разных опытах. Там анализировался факт отключения нагрузки при БАВР, и то, что дело не в скорости работы, колебания и т.д. А в том, в какой момент этого колебания прост случилась авария (отключение выключателя).}{Время (от начала аварии до команд БАВР)}}} \\ \hline
\multicolumn{7}{|c|}{\textbf{Отключения ввода №2 (Т2)}} \\ \hline
\multirow{3}{*}{1} & № работающего ПЧ & 1 & 5 & \multirow{3}{*}{491~А} & \multirow{3}{*}{4,6} & \multirow{3}{*}{31,667~мс} \\ \cline{2-4}
 & Срабатывание БАВР & + & -- & & & \\ \cline{2-4}
 & Уровень колебаний & 5\% & 34\% & & & \\ \hline
\multirow{3}{*}{2} & № работающего ПЧ & 1 & 6 & \multirow{3}{*}{305~А} & \multirow{3}{*}{2,3} & \multirow{3}{*}{33,333~мс} \\ \cline{2-4}
 & Срабатывание БАВР & + & + & & & \\ \cline{2-4}
 & Уровень колебаний & 10\% & 38\% & & & \\ \hline
\multirow{3}{*}{3} & № работающего ПЧ & 3 & 5 & \multirow{3}{*}{320~А} & \multirow{3}{*}{3} & \multirow{3}{*}{34,167~мс} \\ \cline{2-4}
 & С Срабатывание БАВР & + & + & & & \\ \cline{2-4}
 & Уровень колебаний & 8\% & 34\% & & & \\ \hline
\multicolumn{7}{|c|}{\textbf{Отключения ввода №1 (Т1)}} \\ \hline
\multirow{3}{*}{4} & № работающего ПЧ & 1 & 6 & \multirow{3}{*}{270~А} & \multirow{3}{*}{2,3} & \multirow{3}{*}{17,5~мс} \\ \cline{2-4}
 & Срабатывание БАВР & -- & + & & & \\ \cline{2-4}
 & Уровень колебаний & 28\% & 31\% & & & \\ \hline
\multirow{3}{*}{5} & № работающего ПЧ & 2 & 5 & \multirow{3}{*}{350~А} & \multirow{3}{*}{3,2} & \multirow{3}{*}{29,167~мс} \\ \cline{2-4}
 & Срабатывание БАВР & + & + & & & \\ \cline{2-4}
 & Уровень колебаний & 36\% & 22\% & & & \\ \hline
\multirow{3}{*}{6} & № работающего ПЧ & 3 & 5 & \multirow{3}{*}{242~А} & \multirow{3}{*}{2,2} & \multirow{3}{*}{30,833~мс} \\ \cline{2-4}
 & Срабатывание БАВР & + & + & & & \\ \cline{2-4}
 & Уровень колебаний & 17\% & 40\% & & & \\ \hline
\end{tabularx}
\end{table}
